\section{Mixed moments and the remaining case}\label{sec:mixed}

Suppose that a complete set of seven mutually unbiased bases exists and that
its transition matrices satisfy
\begin{equation}\label{eq:all-M}
 g_{H_{rs}}(P\alpha)=0
 \qquad(r\ne s,\;P\text{ a coordinate permutation}).
\end{equation}
We seek a nonnegative averaged quadratic form whose expansion in bounded
moments has a strictly negative upper bound. Corollary~\ref{cor:dichotomy} and
Proposition~\ref{prop:cubic-exclusion} will then exclude a complete set.

\subsection{Sampling and positive forms}

Choose a reference basis index $r$ uniformly from the seven indices. Given
$r$, choose independently a uniform permutation $\sigma$ of its six
coordinates and a uniform bijection $\pi$ from $\{0,\ldots,5\}$ to the
six other basis indices. In this section write
\[
 H_s(i,k)=\sqrt6\,\langle b_{r,\sigma(i)},b_{\pi(s),k}\rangle,
 \qquad g_s=g_{H_s}.
\]
The column order within each basis may remain fixed: each $g_s$ already
averages over all six columns.
For labels $\mathbf s=(s_1,\ldots,s_n)$ and balanced charges
$\mathbf a=(a_1,\ldots,a_n)$ with $\sum_j a_j=0$, define
\begin{equation}\label{eq:mixed-moment}
 \mathcal M(\mathbf s;\mathbf a)
   =\E\,\Ree\!\left(\prod_{j=1}^n g_{s_j}(a_j)\right).
\end{equation}
Equal labels always denote the same matrix. The marginal assignment of any
$l$ distinct labels is a uniform injection into the six nonreference
bases. Thus the definition does not depend on which unused labels are
included in the sampled bijection. The empty product has moment~$1$.

We also use the ordered-pair average of the injective moments of
Section~\ref{sec:pair}:
\begin{equation}\label{eq:average-pair-moment}
 m(A)=\frac1{42}\sum_{r\ne s}m_{H_{rs}}(A),
 \qquad
 H_{rs}(i,k)=\sqrt6\,\langle b_{r,i},b_{s,k}\rangle.
\end{equation}
Here $m_H$ is the real injective moment with the transpose averaging
specified there. All these actual moments obey
\begin{equation}\label{eq:moment-bounds}
 |m(A)|\le1,\qquad |\mathcal M(\mathbf s;\mathbf a)|\le1.
\end{equation}
For the first inequality the integrand has modulus one. For the second,
each $g_s(a)$ is an average of six numbers of modulus one.

The next construction expresses the real parts of Gram entries as linear
combinations of three-factor mixed moments. The third factor comes from
one column shared by all the functions.

In addition to the preceding sampling, choose a uniform column $z$ of
$H_0$. For every balanced charge $a$, define functions on this probability
space by
\begin{equation}\label{eq:F-functions}
 F_0(a)=g_0(a)z^{-a},\qquad
 F_1(a)=\left(\frac15\sum_{s=1}^5g_s(a)\right)z^{-a}.
\end{equation}
The same sampled column $z$ is used in every function $F_t(a)$.
This use of moments and positive Gram forms follows the broader
polynomial-optimization approach to MUBs
\cite{BrierleyWeigert2010SDP,GriblingPolak2024}; the shared-column
functions above specify the particular forms used here.

In the formulas below, an index $1$ on $F$ denotes the average over five
bases in \eqref{eq:F-functions}; a label $1$ in $\mathcal M$ denotes a
representative of those bases after taking expectation.

\begin{lemma}[Gram entries]\label{lem:gram}
For balanced $a,v$, put $\mathbf b=(a,-v,v-a)$. If
$(s,t)\in\{0,1\}^2$ and $(s,t)\ne(1,1)$, then
\begin{equation}\label{eq:gram-first}
 \Ree\,\E\bigl(F_s(a)\overline{F_t(v)}\bigr)
       =\mathcal M((s,t,0);\mathbf b).
\end{equation}
For the remaining case,
\begin{equation}\label{eq:gram-second}
 \Ree\,\E\bigl(F_1(a)\overline{F_1(v)}\bigr)
   =\frac15\mathcal M((1,1,0);\mathbf b)
    +\frac45\mathcal M((1,2,0);\mathbf b).
\end{equation}
\end{lemma}

\begin{proof}
Conditioning first on all the matrices and averaging over $z$ gives
$\E_z z^{v-a}=g_0(v-a)$. In a term with exactly one $F_1$, the five
summands have the same expectation by label exchangeability. In a term
with two $F_1$ factors, expand the two sums independently. Of the $25$
ordered choices of labels from $\{1,\ldots,5\}$, five have equal labels
and twenty have distinct labels. Their marginal distributions give
\eqref{eq:gram-second}. This averaging removes the auxiliary column
completely. A zero charge contributes $g_s(0)=1$ and may be deleted;
removing any unused label preserves the uniform injection law already
described.
\end{proof}

\subsection{Reducing the Gram entries}

The three-factor moments in Lemma~\ref{lem:gram} can be simplified using
symmetry, character vanishing, and the identities of a complete family.
When all factors refer to one matrix, the injective moments and equations
from Section~\ref{sec:pair} apply as well.

\Needspace{7\baselineskip}
\begin{lemma}[Mixed-moment reductions]\label{lem:mixed-reductions}
Under \eqref{eq:all-M}, the following operations and identities are valid.
\begin{enumerate}[label=\textup{(\roman*)}]
 \item A common coordinate permutation of all charges, simultaneous
 negation of all charges, a simultaneous reordering of factors and their
 labels, and an arbitrary injective renaming of the used labels preserve
 \eqref{eq:mixed-moment}. A zero charge factor may be omitted. Coordinates
 on which every charge is zero may also be omitted, with the remaining
 coordinates interpreted as a uniform injection into the six row positions.

 \item The moment vanishes if any charge has order one or is a coordinate
 permutation of $\alpha$.

 \item Suppose that $a_j$ has order two and its label $s_j$ occurs only
 once. Let $S$ be the set of distinct used labels and $l=|S|$. For each
 $t\in S\setminus\{s_j\}$, let $\mathbf s^{j\to t}$ be obtained from
 $\mathbf s$ by replacing $s_j$ with $t$. Then
 \begin{equation}\label{eq:singleton-reduction}
 \mathcal M(\mathbf s;\mathbf a)
 =-\frac1{7-l}\sum_{t\in S\setminus\{s_j\}}
       \mathcal M(\mathbf s^{j\to t};\mathbf a).
 \end{equation}

 \item If all labels are equal, the moment is a linear combination of
 the pair moments \eqref{eq:average-pair-moment}. More precisely, let
 $A$ have columns $a_1,\ldots,a_n$, after removing its zero rows. For a
 set partition $\mathcal P$ of $\{1,\ldots,n\}$, let $A_{\mathcal P}$
 have one column $\sum_{j\in B}a_j$ for each block $B\in\mathcal P$.
 Then
 \begin{equation}\label{eq:single-label-partitions}
 \mathcal M((s,\ldots,s);\mathbf a)
 =\sum_{\substack{\mathcal P\text{ a set partition}\\|\mathcal P|\le6}}
       \frac{(6)_{|\mathcal P|}}{6^n}\,m(A_{\mathcal P}),
 \end{equation}
 where $(6)_k=6\cdot5\cdots(6-k+1)$. Zero rows or columns in
 $A_{\mathcal P}$ may be deleted as in Section~\ref{sec:pair}.

 \item For balanced $a,b,c$ with $a+b+c=0$,
 \begin{align}
 &\sum_{j=0}^5
 \mathcal M\bigl((s,s,u);
       (a+e_i-e_j,b-e_i+e_j,c)\bigr)
   =\mathcal M\bigl((s,u);(a+b,c)\bigr),
       \label{eq:mixed-same}\\
 &\sum_{d\in D}
 \mathcal M\bigl((s,t,u);(a+d,b-d,c)\bigr)=0
       \qquad(s\ne t).
       \label{eq:mixed-different}
 \end{align}
 Equation~\eqref{eq:mixed-same} also holds with the signs of both shifts
 reversed. The third label $u$ may coincide with another label.
\end{enumerate}
\end{lemma}

\begin{proof}
Uniform sampling gives the coordinate and label symmetries in (i), while
commutativity gives factor reordering. Simultaneous negation conjugates
the whole product, leaving its real part unchanged. A zero charge gives
$g_s(0)=1$. Removing unused labels or coordinates leaves a uniform
injection on those retained. Part (ii) follows from Hadamard
orthogonality and \eqref{eq:all-M}. In particular, a permuted $\alpha$
has six nonzero coordinates; deleting zero coordinate rows cannot change
whether a charge is of this type.

For (iii), fix $r$, the row permutation, and the bases assigned to the
$l-1$ other used labels. There are $6-(l-1)=7-l$ available bases for the
singleton label. For this fixed row order, let $\widetilde g_v$ denote
the character average for the external basis $v$. Equation~\eqref{eq:order-two} gives
\[
 \sum_{\substack{v\ne r\\v\text{ not already assigned}}}
       \widetilde g_v(a_j)
 =-\sum_{t\in S\setminus\{s_j\}}g_t(a_j),
\]
with the same row order throughout. Multiply by the other factors and
average over the remaining choices. The marginal distribution of the
$l-1$ retained labels is exactly the distribution defining each moment
on the right of \eqref{eq:singleton-reduction}.

For (iv), expand each $g_s$ as its column average and group the resulting
$n$ column choices by their equality partition. A partition with $k$
blocks has $(6)_k$ distinct assignments to its blocks, out of $6^n$
choices, giving the displayed coefficient. The retained rows are sampled
injectively. It remains to justify the transpose averaging in $m(A)$.
The marginal ordered pair of bases is uniform, and
$H_{sr}=\overline{H_{rs}}^{\,\mathsf T}$. Exchanging $r$ and $s$
exchanges the row and column injections and sends an exponent matrix
$A$ to $-A^{\mathsf T}$. Negating all exponents conjugates the monomial,
leaving its real part unchanged. Thus the ordered-pair average is already
invariant under transposition. This proves \eqref{eq:single-label-partitions}, including
the transpose convention in \eqref{eq:average-pair-moment}.

Finally, multiply \eqref{eq:same-matrix} or
\eqref{eq:different-matrices} by $g_u(c)$ and take real expectations.
These pointwise identities give (v).
\end{proof}

In the finite calculation below, a mixed key consists of its factor
labels and the matrix whose columns are its charges. The reductions in
Lemma~\ref{lem:mixed-reductions} apply to these keys. Single-label keys
are expanded into the moments $m(A)$, and the pair relations from
Section~\ref{sec:pair} may then be used. If a single-label expression
vanishes by part (ii), its expansion also gives a valid homogeneous
relation among the $m(A)$. These steps preserve the interpretation of
every remaining symbol as an actual moment.

\subsection{The six positive forms}

Here are the finite collections of functions and integer weights used
in the certificate. The exact data are the literal list named
\texttt{data} in Appendix~\ref{app:code}, also supplied in the companion
file \texttt{certificate/verify.py}. Each record has the form
\[
 (\lambda,\ \texttt{ws},\ e_\lambda,\ \texttt{a}).
\]
The list $\texttt{ws}=((t_0,v_0),\ldots,(t_{m-1},v_{m-1}))$ specifies
$m$ pairs with $t_i\in\{0,1\}$ and balanced $v_i\in\Z^6$.
The integer list $\texttt{a}$ specifies a symmetric $m\times m$ matrix
$Z_\lambda$ by its lower triangle in row order: with zero-based list
indices,
\begin{equation}\label{eq:Z-encoding}
 (Z_\lambda)_{ij}=(Z_\lambda)_{ji}
       =\texttt{a}_{\,i(i+1)/2+j}
       \qquad(0\le j\le i<m).
\end{equation}
Thus the data define the matrices and lists without any optimization or
choice of numerical tolerances. Their sizes and scales are
\begin{center}
\begin{tabular}{@{}lrr@{}}
\toprule
$\lambda$ & $m$ & $e_\lambda$\\
\midrule
$(1,1,1,2)$ & 5  & 14\\
$(1,1,3)$   & 21 & 13\\
$(1,2,2)$   & 11 & 13\\
$(1,4)$     & 27 & 14\\
$(2,3)$     & 22 & 14\\
$(5)$       & 12 & 15\\
\bottomrule
\end{tabular}
\end{center}

For each of these six lists of row lengths $\lambda$, place the positions
$1,\ldots,5$, in order, in successive left-justified rows of lengths
$\lambda_1,\lambda_2,\ldots$. Let $R_\lambda$ consist of permutations
within these rows, and let $C_\lambda$ consist of permutations within
the resulting columns. Both groups act on $\{0,\ldots,5\}$ by fixing
position~$0$. Our convention for their action on charges is
\[
 (P_pv)_i=v_{p(i)}.
\]
These signed row and column permutation sums are the usual
Young-symmetrizer construction. Symmetry reduction of positive Gram
forms is developed generally in \cite{GatermannParrilo2004} and for MUB
moment problems in \cite{GriblingPolak2024}. Here the finite sums are
given explicitly, so their use requires no representation-theoretic
decomposition theorem.
Define the vector $Y_\lambda$ by
\begin{equation}\label{eq:Young-functions}
 (Y_\lambda)_i
    =\sum_{\substack{p\in R_\lambda\\q\in C_\lambda}}
          \operatorname{sgn}(q)\,
          F_{t_i}(P_pP_qv_i)
       \qquad(0\le i<m).
\end{equation}
In particular, $(P_pP_qv)_i=v_{q(p(i))}$, which is the order implemented
by \texttt{perm(perm(v,q),p)} in the appendix. Formula
\eqref{eq:Young-functions} defines a finite linear combination of the
functions \eqref{eq:F-functions} for every listed pair.

\begin{proposition}[Mixed-moment certificate]\label{prop:mixed-certificate}
Each matrix $Z_\lambda$ specified above is positive definite. Under the
mixed-moment identities of Lemma~\ref{lem:mixed-reductions}, the forms
\eqref{eq:Young-functions} satisfy
\begin{align}
 5\sum_\lambda 10^{15-e_\lambda}
       \E\bigl(Y_\lambda^*Z_\lambda Y_\lambda\bigr)
 &\le
       \frac{-613302797399911+179721388988719}{6480}
       \label{eq:mixed-certificate}\\
 &=-\frac{2007321335237}{30}<0.\notag
\end{align}
\end{proposition}

\begin{proof}
The exact finite calculations establishing the matrix and coefficient
claims are detailed in Section~\ref{sec:verification}, with their full
integer and rational arithmetic in Appendix~\ref{app:code}. We specify
here the mathematical expansion to which those checks apply.

First expand each $Y_\lambda$ by \eqref{eq:Young-functions}, and replace
every real Gram entry using Lemma~\ref{lem:gram}. Multiplication by
five makes its weights either $5$ or the pair $(1,4)$. Because
$Z_\lambda$ is real symmetric, summing its diagonal contributions and
twice its lower off-diagonal contributions gives the real quantity
$\E(Y_\lambda^*Z_\lambda Y_\lambda)$. The factor
$10^{15-e_\lambda}$ is retained for each block.

Apply the mixed reductions and the homogeneous relations
\eqref{eq:mixed-same}--\eqref{eq:mixed-different}, using the finite
selection of relations with total charge order at most $11$ described
in Section~\ref{sec:verification}. That selection only restricts which
valid relations are used. The exact substitutions give
\begin{equation}\label{eq:mixed-residual}
 5\sum_\lambda 10^{15-e_\lambda}
       \E\bigl(Y_\lambda^*Z_\lambda Y_\lambda\bigr)
    = f+\sum_{\rho\in\mathcal R}\beta_\rho X_\rho,
\end{equation}
where each $X_\rho$ is a nonconstant symbol representing one of the
original moments $m(A)$ or $\mathcal M(\mathbf s;\mathbf a)$, and
\begin{equation}\label{eq:mixed-coefficients}
 f=-\frac{613302797399911}{6480},\qquad
 \sum_{\rho\in\mathcal R}|\beta_\rho|
     =\frac{179721388988719}{6480}.
\end{equation}
The linear-algebra verification in Section~\ref{sec:verification}
proves that these substitutions follow from the chosen equations.
The bounds \eqref{eq:moment-bounds} applied to
\eqref{eq:mixed-residual} now prove
\eqref{eq:mixed-certificate}.

For each $Z_\lambda$, the exact Schur-complement calculation in
Section~\ref{sec:verification} verifies positive pivots and hence positive
definiteness.
\end{proof}

\begin{proof}[Completion of the proof of Theorem~\ref{thm:main}]
If seven mutually unbiased bases existed, Corollary~\ref{cor:dichotomy}
would give case (M) or case (C) for each transition matrix.
Proposition~\ref{prop:cubic-exclusion} excludes case (C), so all pairs
satisfy \eqref{eq:all-M} and the preceding construction applies.
For a complex vector $Y=x+\mathrm i y$ and a real symmetric positive
definite matrix $Z$,
\[
 Y^*ZY=x^{\mathsf T}Zx+y^{\mathsf T}Zy\ge0.
\]
Consequently the left side of \eqref{eq:mixed-certificate} is
nonnegative. Its strictly negative upper bound is a contradiction.
This proves Theorem~\ref{thm:main}.
\end{proof}
